\chapter{Security in TUI apps}
\label{ch:security2}

\begin{aside}
TUIs run with the user's privileges. Reading a file the user
can read is fine. Writing one they didn't intend to write is
a bug. Executing one they didn't intend to execute is a
disaster.
\end{aside}

\section{Input validation}

User input is untrusted. A filename from a search box can
contain \code{../}. A URL can be malicious. Sanitise.

\section{Shell-out hygiene}

If your app shells out (e.g., \code{system("foo " +
arg)}), the arg is a shell injection vector. Use
\code{execvp} with separate arguments, never string
concatenation.

\section{Path traversal}

Joining user input to a base path can escape via
\code{../}. Resolve and verify:

\begin{lstlisting}
auto p = std::filesystem::canonical(base / user_input);
if (!p.starts_with(base)) throw "out of bounds";
\end{lstlisting}

\section{Secret handling}

API tokens, passwords. Never write to log files. Never
display in plain text. Use \code{mlock} on Linux to keep
out of swap. Wipe with \code{explicit\_bzero} on release.

\section{Code injection in displayed strings}

If you display a string with embedded ANSI escapes from
untrusted source, the user's terminal interprets them ---
potentially moving the cursor, setting the title to
something nasty, or worse.

\maya{} provides \code{sanitise\_for\_display} that strips
control bytes from untrusted strings before they reach
the cell buffer.

\section{Update mechanism}

If your app self-updates, the update channel is a security
boundary. Sign updates; verify on download. \maya{}
doesn't ship an updater; if you write one, follow
established practice.

\section{TLS}

Network connections to anything sensitive must use TLS.
Verify certificates; never disable verification ``for
testing''.

\section{Recap}

\begin{recap}
\begin{itemize}[leftmargin=*]
  \item Validate user input.
  \item Never shell concatenate.
  \item Canonicalise paths.
  \item Secrets out of logs, out of swap, wiped on release.
  \item Sanitise untrusted strings before display.
  \item Signed updates over TLS.
\end{itemize}
\end{recap}

\section*{Workshop}

\begin{enumerate}[leftmargin=*]
  \item Audit your app's shell-outs; replace with
    \code{execvp}.
  \item Add path-traversal protection to file pickers.
  \item Strip control bytes from any string from
    untrusted sources before display.
\end{enumerate}
